% Proof package P3 (result R10). Compile from the folder proofs/:
%   latexmk -pdf -interaction=nonstopmode -halt-on-error -outdir=build P3_research.tex
% The registry macros come from ../scripts/P3_research/numbers_P3.tex (written by P3_research.py).
\documentclass[12pt]{article}
\usepackage[margin=1in]{geometry}
\usepackage{amsmath,amssymb,amsthm,booktabs}
\usepackage{newtxtext,newtxmath}
\usepackage{enumitem}
\usepackage[hidelinks]{hyperref}

\newtheorem{proposition}{Proposition}
\newtheorem{lemma}{Lemma}
\newtheorem{corollary}{Corollary}
\theoremstyle{definition}
\newtheorem{definition}{Definition}
\theoremstyle{remark}
\newtheorem{remark}{Remark}

\newcommand{\E}{\mathbb E}
\newcommand{\Ent}{\mathsf E}
\newcommand{\Own}{\mathsf O}
\newcommand{\fr}{\mathfrak r}
\newcommand{\DT}{\Delta_T}
\newcommand{\lh}{\hat\lambda}
\newcommand{\lb}{\bar\lambda}
\newcommand{\st}[1]{\textsc{#1}}

% Registry numbers: \reg{<key>}, generated by scripts/P3_research/P3_research.py.
% numbers_P3.tex: generated by scripts/P3_research/P3_research.py from registry/P3_research.csv.
% Do not edit by hand.
\makeatletter
\providecommand{\reg}[1]{\ifcsname reg@#1\endcsname\csname reg@#1\endcsname\else\textbf{??\detokenize{#1}}\fi}
\expandafter\def\csname reg@res_a\endcsname{1/40}
\expandafter\def\csname reg@res_cap\endcsname{0.336122}
\expandafter\def\csname reg@res_lambda_r155\endcsname{0.237901}
\expandafter\def\csname reg@res_entry_r155\endcsname{0.507630}
\expandafter\def\csname reg@res_muplus_r155\endcsname{0.673142}
\expandafter\def\csname reg@res_own_r155\endcsname{0.298422}
\expandafter\def\csname reg@res_lamE_r155\endcsname{0.196741}
\expandafter\def\csname reg@res_Delta_r155\endcsname{0.097581}
\expandafter\def\csname reg@res_tau_r155\endcsname{0.646798}
\expandafter\def\csname reg@res_lamgap_r155\endcsname{0.041160}
\expandafter\def\csname reg@res_capgap_r155\endcsname{0.098221}
\expandafter\def\csname reg@res_UH_r155\endcsname{0.011895}
\expandafter\def\csname reg@res_ULgap_r155\endcsname{0.007802}
\expandafter\def\csname reg@res_lambda_r160\endcsname{0.279037}
\expandafter\def\csname reg@res_entry_r160\endcsname{0.513965}
\expandafter\def\csname reg@res_muplus_r160\endcsname{0.698206}
\expandafter\def\csname reg@res_own_r160\endcsname{0.309302}
\expandafter\def\csname reg@res_lamE_r160\endcsname{0.199698}
\expandafter\def\csname reg@res_Delta_r160\endcsname{0.112500}
\expandafter\def\csname reg@res_tau_r160\endcsname{0.648734}
\expandafter\def\csname reg@res_lamgap_r160\endcsname{0.079340}
\expandafter\def\csname reg@res_capgap_r160\endcsname{0.057085}
\expandafter\def\csname reg@res_UH_r160\endcsname{0.013952}
\expandafter\def\csname reg@res_ULgap_r160\endcsname{0.005938}
\expandafter\def\csname reg@res_lambda_r165\endcsname{0.316683}
\expandafter\def\csname reg@res_entry_r165\endcsname{0.519763}
\expandafter\def\csname reg@res_muplus_r165\endcsname{0.720112}
\expandafter\def\csname reg@res_own_r165\endcsname{0.319260}
\expandafter\def\csname reg@res_lamE_r165\endcsname{0.202664}
\expandafter\def\csname reg@res_Delta_r165\endcsname{0.128030}
\expandafter\def\csname reg@res_tau_r165\endcsname{0.650670}
\expandafter\def\csname reg@res_lamgap_r165\endcsname{0.114019}
\expandafter\def\csname reg@res_capgap_r165\endcsname{0.019439}
\expandafter\def\csname reg@res_UH_r165\endcsname{0.015834}
\expandafter\def\csname reg@res_ULgap_r165\endcsname{0.003996}
\expandafter\def\csname reg@res_th_estar_r155\endcsname{0.865703}
\expandafter\def\csname reg@res_th_entry_r155\endcsname{0.456294}
\expandafter\def\csname reg@res_th_estar_r160\endcsname{0.758777}
\expandafter\def\csname reg@res_th_entry_r160\endcsname{0.420777}
\expandafter\def\csname reg@res_th_estar_r165\endcsname{0.673748}
\expandafter\def\csname reg@res_th_entry_r165\endcsname{0.392494}
\expandafter\def\csname reg@res_rR\endcsname{1.501023}
\expandafter\def\csname reg@res_rcap\endcsname{1.677633}
\expandafter\def\csname reg@res_entry_rR\endcsname{0.500846}
\expandafter\def\csname reg@res_lambda_rR\endcsname{0.193853}
\expandafter\def\csname reg@res_entry_cap\endcsname{0.522757}
\expandafter\def\csname reg@res_amono\endcsname{0.161713}
\expandafter\def\csname reg@res_amono_third\endcsname{0.150885}
\expandafter\def\csname reg@res_phi_r155\endcsname{0.004629}
\expandafter\def\csname reg@res_phip_lb\endcsname{0.033767}
\expandafter\def\csname reg@res_rA\endcsname{1.393958}
\expandafter\def\csname reg@res_dphi_lead\endcsname{0.039942}
\expandafter\def\csname reg@res_bound_dDelta\endcsname{11/72}
\expandafter\def\csname reg@res_bound_tau74\endcsname{19/55}
\expandafter\def\csname reg@res_bound_Delta74\endcsname{9/56}
\expandafter\def\csname reg@res_bound_taup\endcsname{16875/211288}
\expandafter\def\csname reg@res_lamE_r175\endcsname{0.208635}
\expandafter\def\csname reg@res_rN\endcsname{1.747878}
\expandafter\def\csname reg@res_rC\endcsname{3.592659}
\expandafter\def\csname reg@margin_res_bound_r1_B\endcsname{0.016421}
\expandafter\def\csname reg@margin_res_tau_r1_B\endcsname{0.026059}
\expandafter\def\csname reg@margin_res_bound_r2_B\endcsname{0.034723}
\expandafter\def\csname reg@margin_res_tau_r2_B\endcsname{0.000334}
\expandafter\def\csname reg@res_entry_r1\endcsname{0.522757}
\expandafter\def\csname reg@res_entry_r2\endcsname{0.250000}
\expandafter\def\csname reg@res_rUR\endcsname{2.434046}
\expandafter\def\csname reg@res_region_cap_0\endcsname{0.290539}
\expandafter\def\csname reg@res_region_alo_0\endcsname{0.030579}
\expandafter\def\csname reg@res_region_aU3_0\endcsname{0.100388}
\expandafter\def\csname reg@res_region_a36_0\endcsname{0.169489}
\expandafter\def\csname reg@res_region_width_0\endcsname{0.006184}
\expandafter\def\csname reg@res_region_cap_1\endcsname{0.300000}
\expandafter\def\csname reg@res_region_alo_1\endcsname{0.028437}
\expandafter\def\csname reg@res_region_aU3_1\endcsname{0.094155}
\expandafter\def\csname reg@res_region_a36_1\endcsname{0.159824}
\expandafter\def\csname reg@res_region_width_1\endcsname{0.008327}
\expandafter\def\csname reg@res_region_cap_2\endcsname{0.310000}
\expandafter\def\csname reg@res_region_alo_2\endcsname{0.026328}
\expandafter\def\csname reg@res_region_aU3_2\endcsname{0.088015}
\expandafter\def\csname reg@res_region_a36_2\endcsname{0.150299}
\expandafter\def\csname reg@res_region_width_2\endcsname{0.010436}
\expandafter\def\csname reg@res_region_cap_3\endcsname{0.320000}
\expandafter\def\csname reg@res_region_alo_3\endcsname{0.024364}
\expandafter\def\csname reg@res_region_aU3_3\endcsname{0.082293}
\expandafter\def\csname reg@res_region_a36_3\endcsname{0.141417}
\expandafter\def\csname reg@res_region_width_3\endcsname{0.012400}
\expandafter\def\csname reg@res_region_cap_4\endcsname{0.330000}
\expandafter\def\csname reg@res_region_alo_4\endcsname{0.022531}
\expandafter\def\csname reg@res_region_aU3_4\endcsname{0.076951}
\expandafter\def\csname reg@res_region_a36_4\endcsname{0.133120}
\expandafter\def\csname reg@res_region_width_4\endcsname{0.014232}
\expandafter\def\csname reg@res_region_cap_5\endcsname{0.336122}
\expandafter\def\csname reg@res_region_alo_5\endcsname{0.021469}
\expandafter\def\csname reg@res_region_aU3_5\endcsname{0.073853}
\expandafter\def\csname reg@res_region_a36_5\endcsname{0.128306}
\expandafter\def\csname reg@res_region_width_5\endcsname{0.015294}
\expandafter\def\csname reg@res_region_cap_6\endcsname{0.336721}
\expandafter\def\csname reg@res_region_alo_6\endcsname{0.021368}
\expandafter\def\csname reg@res_region_aU3_6\endcsname{0.073556}
\expandafter\def\csname reg@res_region_a36_6\endcsname{0.127845}
\expandafter\def\csname reg@res_region_width_6\endcsname{0.015396}
\expandafter\def\csname reg@res_cap_lo\endcsname{0.290539}
\expandafter\def\csname reg@res_cap_hi\endcsname{0.336721}
\expandafter\def\csname reg@res_a_up\endcsname{0.036763}
\expandafter\def\csname reg@res_a_lo\endcsname{0.021469}
\expandafter\def\csname reg@res_cap_lo_a\endcsname{0.316683}
\expandafter\def\csname reg@margin_res_aup_B\endcsname{0.011763}
\expandafter\def\csname reg@margin_res_alo_B\endcsname{0.003531}
\expandafter\def\csname reg@margin_res_caplo_B\endcsname{0.045583}
\expandafter\def\csname reg@margin_res_caphi_B\endcsname{0.000599}
\expandafter\def\csname reg@res_a_floor3\endcsname{0.100388}
\expandafter\def\csname reg@res_aU3_B\endcsname{0.073853}
\expandafter\def\csname reg@res_beta_lo\endcsname{0.236157}
\expandafter\def\csname reg@res_beta_hi\endcsname{0.4897}
\expandafter\def\csname reg@res_R_access_r1\endcsname{0.872531}
\expandafter\def\csname reg@res_R_hidden_r1\endcsname{0.614583}
\expandafter\def\csname reg@res_dR_r1\endcsname{0.257948}
\expandafter\def\csname reg@res_W_access_r1\endcsname{2.634044}
\expandafter\def\csname reg@res_W_hidden_r1\endcsname{2.552083}
\expandafter\def\csname reg@res_dW_r1\endcsname{0.081960}
\expandafter\def\csname reg@res_cost_r1\endcsname{0.001412}
\expandafter\def\csname reg@res_tcost_r1\endcsname{0.006722}
\expandafter\def\csname reg@res_calA_lambda_r155\endcsname{0.250249}
\expandafter\def\csname reg@res_calA_entry_r155\endcsname{0.531281}
\expandafter\def\csname reg@res_calA_lambda_r160\endcsname{0.293325}
\expandafter\def\csname reg@res_calA_entry_r160\endcsname{0.536666}
\expandafter\def\csname reg@res_calA_lambda_r165\endcsname{0.332615}
\expandafter\def\csname reg@res_calA_entry_r165\endcsname{0.541577}
\expandafter\def\csname reg@res_Delta_r0\endcsname{1/60}
\expandafter\def\csname reg@margin_res_weak_B\endcsname{0.003333}
\expandafter\def\csname reg@margin_res_nosale_B\endcsname{0.007992}
\expandafter\def\csname reg@margin_res_half_B\endcsname{1.195833}
\expandafter\def\csname reg@margin_res_beta_B\endcsname{0.038678}
\expandafter\def\csname reg@margin_res_rR_B\endcsname{0.048977}
\expandafter\def\csname reg@margin_res_rcap_B\endcsname{0.027633}
\expandafter\def\csname reg@margin_res_rN_B\endcsname{0.070245}
\expandafter\def\csname reg@margin_res_amono_B\endcsname{0.136713}
\expandafter\def\csname reg@margin_res_lamcap_B\endcsname{0.019439}
\makeatother


\title{Competition creates research\\
\large Proof package P3 (result R10): the order slot with costly research, general $\beta$ and calibration B}
\author{}
\date{23 October 2026}

\begin{document}
\maketitle

\section*{Summary}

This package proves result R10 for the order slot. Before it learns the class, the investor chooses a private
research intensity. Research succeeds with that probability. Success reveals the class, and the investor then
fills the order slot. After a failure, a liquidity trader fills the slot. The research intensity is therefore the
presence probability $\lambda$ of the fixed-$\lambda$ model, chosen at a quadratic cost and bounded by a cap
$\lb$. The cap equals the benchmark presence probability $\lambda_B$, so research at the cap reproduces the
fixed-$\lambda$ model at calibration B.

The results, with their status labels:
\begin{enumerate}[label=(\arabic*),nosep]
\item The equilibrium object and the joint deviation in research and orders (Definition~\ref{def:eq},
  Lemma~\ref{lem:joint}). \st{analytical}.
\item At $r_0=6/5$, zero research is the unique outcome (Proposition~\ref{prop:weak}). \st{analytical}.
\item At $r\in\{1.55,1.60,1.65\}$, the complete equilibrium set: zero research and a research branch under the
  challenger rule C1, plus one threshold equilibrium under rule CM (Proposition~\ref{prop:middle}).
  \st{analytical}; the values at B are \st{computer-assisted} enclosures.
\item Along the research branch, research and entry rise strictly with incumbent strength. At B the branch
  occupies $[r_R,\bar r_R)=[\reg{res_rR},\reg{res_rcap})$ and coexists with zero research
  (Proposition~\ref{prop:branch}). \st{analytical} at general $\beta$; the interval at B is
  \st{computer-assisted}.
\item At $r_1=3$ and $r_2=18/5$, research at the cap is the unique outcome, and entry equals the fixed-$\lambda$
  entry (Proposition~\ref{prop:strong}). \st{analytical}, with interval margins.
\item The admissible region at B: every claim holds with strict margins exactly when
  $\lambda_E(3)<\lb<\lambda_E(18/5)$ and $a_{\rm lo}(\lb)<a<a_{\rm up}$ (Proposition~\ref{prop:region}).
  \st{computer-assisted}.
\item Price access at $r_1$ raises proceeds by \reg{res_dR_r1} and net surplus by \reg{res_dW_r1}
  (Proposition~\ref{prop:access}). \st{analytical}, with interval values.
\end{enumerate}
The research cost $a=\reg{res_a}$ of the earlier study works at calibration B. It lies inside the cost window
$(\reg{res_a_lo},\reg{res_a_up})$ at the cap $\lb=\lambda_B=\reg{res_cap}$.

\section{The model with research}\label{sec:model}

\subsection{Primitives}

The auction stage is that of Sections 2.1--2.2 and 3.1 of the manuscript. The classes are point masses $h$ and
$\ell$ with prior $\pi=1/2$. The incumbent value is $R\sim U[0,r]$ with $0<p<\ell<r<h$. The low preparation cost
is $c_L=0$, and the high cost $c_H$ has probability $1-\rho$. The objects $t_0,t_H,t_L,g_H,g_L$,
$\DT=t_H-t_L=(r-\ell)^2/(2r)$, $B_r(\mu)=g_L+\mu(g_H-g_L)$ and $\tau(r)=(c_H-g_L)/(g_H-g_L)$ are those of
equation (4) of the Laplace manuscript. The program \texttt{P3\_research.py} checks each closed form against
direct integration over $R$ at the six strengths.

The slot has the liquidity law
\[
\eta(-1)=\beta,\qquad \eta(0)=1-2\beta,\qquad \eta(1)=\beta,\qquad 0<\beta\le 1/3 .
\]
The research intensity is $\lambda\in[0,\lb]$ with a cap $0<\lb<1$. Its cost is $(a/2)\lambda^2$ with $a>0$.
This package keeps the prior at $\pi=1/2$.

\paragraph{Calibration B.} $h=10$, $\ell=1$, $p=1/2$, $\rho=1/4$, $c_H=6$, $k=1/50$, $\pi=1/2$,
$1+u=e$, $\beta=0.294655$, the cap $\lb=\lambda_B=u\beta/(1+u\beta)=\reg{res_cap}$, and the research cost
$a=\reg{res_a}$. The strengths are $r_0=6/5$, the intermediate strengths $31/20$, $8/5$, $33/20$, and
$r_1=3$, $r_2=18/5$. Every statement that depends on $e$ uses rational bounds on $e$ (Section~\ref{sec:cert}).

\subsection{Timing}

\begin{enumerate}[nosep]
\item Nature draws the class $j$, the incumbent value $R$ and the challenger's cost $C$. The draws are independent.
\item The investor chooses a research intensity $z\in[0,\lb]$ and pays $(a/2)z^2$. Market makers and the challenger do not observe $z$.
\item Research succeeds with probability $z$, independently of $(j,R,C)$. After a success, the investor learns $j$
  and fills the slot with an order $q\in\{-1,0,1\}$ at cost $k|q|$. After a failure, a liquidity trader fills the
  slot with an order drawn from $\eta$. The investor then does not trade.
\item Competitive market makers observe the order $q$ and set $P=\E[V_T\mid q]$. They use the conjectured
  intensity $\lh$ and the conjectured order distributions.
\item The challenger observes the price and its cost, then decides whether to prepare.
\item Prepared bidders bid their values in the second-price auction with reserve $p$.
\end{enumerate}
An order comes from the informed investor with probability $\lambda$, as in the fixed-$\lambda$ slot. Research at
the cap therefore gives the order law of the fixed-$\lambda$ model with presence $\lb$. More research also removes
liquidity mass, since the liquidity trader fills the slot only after a failure.

\subsection{The objects of a candidate}

Fix a conjectured intensity $\lh\in[0,\lb]$ and order distributions $\sigma_H,\sigma_L$ on $\{-1,0,1\}$ after a
success. For each order $x\in\{-1,0,1\}$, the probability of order $x$ in class $j$ and the posterior are
\begin{equation}\label{eq:masses}
f_j(x)=(1-\lh)\eta(x)+\lh\sigma_j(x),\qquad \mu(x)=\frac{f_H(x)}{f_H(x)+f_L(x)} .
\end{equation}
Write $\hat b=(1-\lh)\beta$ and $\hat b_0=(1-\lh)(1-2\beta)$ for the liquidity masses at $\pm1$ and at $0$. The
notation follows the shared notation of the documents: $f_j(x)$, $\mu(x)$, $P(x)$ and $e(x)$ are functions of the
order $x$. The core package writes $a_j(x)$ for $f_j(x)$; the letter $a$ here is the research cost. A proof that
names the investor's order uses $q$ as the argument, as in $U_j(q)$ and $P(q)$.
Let $\zeta_x\in[0,1]$ be the preparation probability of the high-cost challenger after the price of order $x$, and
$e(x)=\rho+(1-\rho)\zeta_x$ total entry. Since $(j,R)$ is independent of $C$,
$\E[V_T\mid j,x]=t_0+e(x)(t_j-t_0)$. Competitive pricing gives
\begin{equation}\label{eq:price}
P(x)=t_0+e(x)\left(t_L-t_0+\DT\mu(x)\right).
\end{equation}
The payoff of order $q$ after a success is $q(\E[V_T\mid j,q]-P(q))-k|q|$, so
\begin{equation}\label{eq:payoffs}
U_H(q)=q\,e(q)\DT(1-\mu(q))-k|q|,\qquad U_L(q)=-q\,e(q)\DT\mu(q)-k|q| .
\end{equation}
The research value is the expected optimal trading payoff after a success:
\begin{equation}\label{eq:V}
V=\tfrac12\max_q U_H(q)+\tfrac12\max_q U_L(q)\ \ge 0 .
\end{equation}
Entry and high-class ownership are
\begin{equation}\label{eq:entry}
\Ent=\sum_q \tfrac12\left[f_H(q)+f_L(q)\right]e(q),\qquad \Own_H=\tfrac12\sum_q f_H(q)e(q) .
\end{equation}
Ownership uses $W_H=1$, since $h>r$.

\paragraph{Derived functions.} For $\lambda\in[0,1)$ put $u_\lambda=\lambda/[(1-\lambda)\beta]$,
\[
M_\lambda=\frac{1+u_\lambda}{2+u_\lambda},\qquad
m_\lambda=1-M_\lambda=\frac{(1-\lambda)\beta}{2\beta+\lambda(1-2\beta)} .
\]
$M_\lambda$ is the bound $M(1/2)$ of Section 3.2 at presence $\lambda$. The function $m_\lambda$ falls strictly from
$1/2$ at $\lambda=0$ to $0$ at $\lambda=1$, because $1-2\beta>0$. At the cap of calibration B,
$u_{\lambda_B}=u$, so $M_{\lambda_B}=e/(1+e)$ and $m_{\lambda_B}=1/(1+e)$.

For $1/2<\tau(r)<1$, the \emph{recruitment intensity} $\lambda_E(r)$ is the unique $\lambda$ with
$M_\lambda=\tau(r)$:
\begin{equation}\label{eq:lamE}
\lambda_E(r)=\frac{\beta\,[2\tau(r)-1]}{\beta+(1-2\beta)[1-\tau(r)]},\qquad
\frac{d\lambda_E}{d\tau}=\frac{\beta}{\{\beta+(1-2\beta)(1-\tau)\}^2}>0 .
\end{equation}
The \emph{research polynomial} is
\begin{equation}\label{eq:F}
F_r(\lambda)=2a(1-2\beta)\lambda^2+\left[4a\beta+k(1-2\beta)+\beta\DT(r)\right]\lambda+\beta\left[2k-\DT(r)\right]
=2\left[2\beta+\lambda(1-2\beta)\right]\left[a\lambda-\tfrac12\left(\DT(r)m_\lambda-k\right)\right].
\end{equation}
All coefficients except the constant are positive, so $F_r$ rises strictly on $\lambda\ge0$. If $\DT(r)>2k$, then
$F_r(0)<0<F_r(1)$, and $F_r$ has a unique root $\lambda_R(r)\in(0,1)$. The \emph{recruitment margin} and the
\emph{threshold entry} are
\begin{equation}\label{eq:phi}
\varphi_a(r)=\DT(r)\left[1-\tau(r)\right]-k-2a\lambda_E(r),\qquad
e^*(r)=\frac{2a\lambda_E(r)+k}{\DT(r)[1-\tau(r)]} .
\end{equation}
Since $m_{\lambda_E}=1-\tau$, equation~\eqref{eq:F} gives
$F_r(\lambda_E(r))=-\left[2\beta+\lambda_E(1-2\beta)\right]\varphi_a(r)$. Thus $\lambda_R(r)\ge\lambda_E(r)$
exactly when $\varphi_a(r)\ge0$, and $e^*(r)<1$ exactly when $\varphi_a(r)>0$.

\subsection{Equilibrium}

\begin{definition}[Equilibrium with research]\label{def:eq}
An equilibrium is a tuple $\mathcal E=(\lh,\sigma_H,\sigma_L,P,\mu,\zeta,\text{bids})$. Its components are:
\begin{itemize}[nosep]
\item $\lh\in[0,\lb]$, the conjectured research intensity;
\item $\sigma_j\in\Delta(\{-1,0,1\})$, the order distribution of class $j$ after a success;
\item $P=(P(x))$, a price on every order, and $\mu=(\mu(x))$, a posterior on every order;
\item a preparation rule: the zero-cost challenger prepares at every price, and the high-cost challenger prepares
  with probability $\zeta(\bar P)$ at each price $\bar P$ in the image of $P$; write $\zeta_x=\zeta(P(x))$;
\item bids: each prepared bidder bids its value.
\end{itemize}
The tuple satisfies conditions (E1)--(E6) of the fixed-$\lambda$ model at the intensity $\lh$, and a seventh
condition for research.
\begin{enumerate}[label=(E\arabic*),nosep]
\item \emph{Bayes consistency.} $\mu(q)$ satisfies \eqref{eq:masses} for every $q$.
\item \emph{Challenger beliefs.} At each price, the challenger's belief is the probability of class $H$ given
  that price, from the joint law of the class and the order.
\item \emph{Entry optimality.} $\zeta_q=1$ if the belief at $P(q)$ gives $B_r>c_H$, and $\zeta_q=0$ if it gives
  $B_r<c_H$. At equality the challenger tie rule applies.
\item \emph{Competitive prices.} $P(q)=\E[V_T\mid q]$ for every $q$, that is \eqref{eq:price}.
\item \emph{Investor optimality.} Against the fixed $P$ and $e$, each $\sigma_j$ puts probability one on the
  maximizers of $U_j$. At indifference the investor tie rule applies.
\item \emph{Bids.} Truthful bids are weakly dominant in the second-price auction.
\item \emph{Research optimality.} Against the fixed $P$ and $e$, $\lh$ is the research choice of a best response
  among all joint deviations in research and orders.
\end{enumerate}
\end{definition}

The conjectured intensity enters (E1)--(E4) through the order masses. Condition (E7) requires the investor's
best response to the price function and entry rule of the tuple to equal that conjecture. By
Lemma~\ref{lem:bounds}(b), (E2) makes the challenger's belief at $P(q)$ equal to $\mu(q)$. Condition (E5) applies to
both classes also when $\lh=0$: a deviation in research makes a success possible, and the order after it must be
optimal.

\paragraph{The joint deviation.} A joint deviation $(z,\sigma'_H,\sigma'_L)$ changes the true order law to
$(1-z)\eta+z\sigma'_j$. It does not change the prices or the entry rule, because market makers and the challenger
use the candidate's $\lh$ and $\sigma_j$. The investor trades only after a success. Its expected payoff is
\begin{equation}\label{eq:Pi}
\Pi(z,\sigma')=z\left[\tfrac12\sum_q\sigma'_H(q)U_H(q)+\tfrac12\sum_q\sigma'_L(q)U_L(q)\right]-\frac a2z^2 ,
\end{equation}
with $U_j$ from \eqref{eq:payoffs} at the candidate's prices and entry.

\paragraph{Tie rules.}
\begin{itemize}[nosep]
\item Investor. IM: any mixture over the orders with the maximal payoff. I0: the zero order whenever it maximizes.
\item Challenger. C1: the high-cost challenger prepares at indifference (the rule of Section 2.2).
  CM: any preparation probability at indifference.
\end{itemize}
Research itself has no tie (Lemma~\ref{lem:joint}).

\section{Lemmas}\label{sec:lemmas}

\begin{lemma}[Joint deviation; \st{analytical}]\label{lem:joint}
For every candidate, $\sup_{\sigma'}\Pi(z,\sigma')=zV-(a/2)z^2$ for each $z\in[0,\lb]$. The unique optimal
research choice is $z^*=\min\{\lb,V/a\}$. Hence (E7) holds if and only if $\lh=\min\{\lb,V/a\}$, and (E5) and (E7)
together give optimality against every joint deviation.
\end{lemma}

\begin{proof}
Fix $z\ge0$. The bracket in \eqref{eq:Pi} is at most $V$, with equality exactly when each $\sigma'_j$ puts mass
only on maximizers of $U_j$. So $\sup_{\sigma'}\Pi(z,\sigma')=zV-(a/2)z^2$. This function of $z$ is strictly
concave, with derivative $V-az$. If $V\ge a\lb$, the derivative is positive on $[0,\lb)$, and the maximizer is
$\lb$. If $V<a\lb$, the maximizer is $V/a\in[0,\lb)$. Every maximizer of $\Pi$ therefore uses $z^*$. For $z^*>0$
it also uses orders that satisfy (E5). The research choice is unique, so no tie rule applies to it.
\end{proof}

\begin{lemma}[Bounds under every profile; \st{analytical}]\label{lem:bounds}
Let $0<\beta\le1/3$ and $0\le\lh\le\lb<1$. For all order distributions $\sigma_H,\sigma_L$ and all
$\zeta\in[0,1]^3$:
\begin{enumerate}[label=(\alph*),nosep]
\item every posterior satisfies $m_{\lh}\le\mu(q)\le M_{\lh}$, and $M_{\lh}\le M_{\lb}$;
\item if (E4) holds and $P(q)=P(q')$, then $\mu(q)=\mu(q')$; so the price reveals the market posterior;
\item $U_H(-1)\le-k$ and $U_L(1)\le-k$; so every equilibrium has $\sigma_H(-1)=\sigma_L(1)=0$;
\item if $\sigma_H(-1)=\sigma_L(1)=0$, then $\mu(-1)\le1/2\le\mu(1)$, and $\mu(1)=M_{\lh}$ when $\sigma_H(1)=1$.
\end{enumerate}
\end{lemma}

\begin{proof}
(a) For $q=\pm1$, $\hat b\le f_j(q)\le\hat b+\lh$. For $q=0$, $\hat b_0\le f_j(0)\le\hat b_0+\lh$, and
$\hat b_0\ge\hat b$ because $\beta\le1/3$. This is the only use of $\beta\le1/3$. Hence
$f_H(q)/f_L(q)\in[1/(1+u_{\lh}),1+u_{\lh}]$ for every $q$, and Bayes' rule at prior $1/2$ gives
$\mu(q)\in[1/(2+u_{\lh}),(1+u_{\lh})/(2+u_{\lh})]=[m_{\lh},M_{\lh}]$. The map $\lambda\mapsto u_\lambda$ rises,
so $M_{\lh}\le M_{\lb}$.

(b) Entry after a price depends only on the price. If $P(q)=P(q')$, then $e(q)=e(q')=:e\ge\rho>0$, and
\eqref{eq:price} gives $e\DT(\mu(q)-\mu(q'))=0$. Since $\DT>0$ for $r>\ell$, $\mu(q)=\mu(q')$. The challenger's
belief after a price is the posterior over the orders with that price, and it equals their common $\mu(q)$.

(c) $U_H(-1)=-e(-1)\DT(1-\mu(-1))-k\le-k<0=U_H(0)$, and $U_L(1)=-e(1)\DT\mu(1)-k\le-k<0=U_L(0)$.

(d) $f_H(1)=\hat b+\lh\sigma_H(1)\ge\hat b=f_L(1)$ and $f_H(-1)=\hat b\le\hat b+\lh\sigma_L(-1)=f_L(-1)$.
With $\sigma_H(1)=1$, $\mu(1)=(\hat b+\lh)/(2\hat b+\lh)=M_{\lh}$.
\end{proof}

\begin{lemma}[Positive research needs a strict full purchase; \st{analytical}]\label{lem:buy}
Let $1/2<\tau(r)$ and consider an equilibrium with $\lh>0$. Then $\sigma_H(1)=1$ and $U_H(1)>0$. Moreover
$\mu(1)=M_{\lh}$, $\mu(0)\le1/2$, $\mu(-1)\le1/2$, and $\zeta_0=\zeta_{-1}=0$.
\end{lemma}

\begin{proof}
By Lemma~\ref{lem:joint}, $\lh=\min\{\lb,V/a\}>0$, so $V>0$. By Lemma~\ref{lem:bounds}(c) and $U_j(0)=0$, either
$U_H(1)>0$ or $U_L(-1)>0$.

Suppose $U_L(-1)>0$. Then $\sigma_L(-1)=1$ by (E5). With Lemma~\ref{lem:bounds}(c), $f_H(-1)=\hat b$ and
$f_L(-1)=\hat b+\lh$, so $\mu(-1)=m_{\lh}\le1/2<\tau(r)$. Hence $\zeta_{-1}=0$ and
$U_L(-1)=\rho\DT m_{\lh}-k>0$. The purchase earns
$U_H(1)=e(1)\DT(1-\mu(1))-k\ge\rho\DT(1-M_{\lh})-k=\rho\DT m_{\lh}-k>0$, by Lemma~\ref{lem:bounds}(a).

So $U_H(1)>0$ in both cases. The purchase is the unique maximizer of $U_H$, and $\sigma_H(1)=1$ under both investor
rules. Lemma~\ref{lem:bounds}(d) gives $\mu(1)=M_{\lh}$ and $\mu(-1)\le1/2$. Also $f_H(0)=\hat b_0\le\hat b_0+\lh\sigma_L(0)
=f_L(0)$, so $\mu(0)\le1/2$. Since $g_H>g_L$, $B_r(\mu)<c_H$ exactly when $\mu<\tau(r)$. So $\tau(r)>1/2$ gives
$\zeta_0=\zeta_{-1}=0$ under both challenger rules.
\end{proof}

At calibration B, $B_r(1/2)=5-r/4+1/(8r)<6=c_H$ for $1<r<10$. So $\tau(r)>1/2$ at every strength of this
package. Also $g_H(r)>c_H$, that is $\tau(r)<1$, for $r\le 18/5$.

\section{Results}\label{sec:results}

\begin{proposition}[Weak incumbent; \st{analytical}]\label{prop:weak}
Let $B_r(1/2)<c_H$, that is, $\tau(r)>1/2$: the challenger with the high cost does not prepare at the prior.
If $\DT(r)<k$, then under every pair of tie rules the unique equilibrium is
$\lh=0$, $\sigma_H=\sigma_L=\delta_0$, $\mu(q)=1/2$, $\zeta_q=0$, $e(q)=\rho$, $P(q)=t_0+\rho(t_L-t_0+\DT/2)$, with
entry $\rho$. At calibration B this holds at $r_0=6/5$, where $\DT=\reg{res_Delta_r0}<k=1/50$ and
$c_H-B_{6/5}(1/2)=\reg{margin_res_half_B}$.
\end{proposition}

The prior-entry hypothesis sets the entry component. Without it, $\DT(r)<k$ still gives zero research and no
trade. But $B_r(1/2)>c_H$ makes both cost types prepare after every order, so entry is $1$. The no-research
condition $\rho\DT\le2k$ of Proposition~\ref{prop:middle} also uses floor entry at the uninformative price. That
proposition carries the same hypothesis through $\tau(r)>1/2$.

\begin{proposition}[Intermediate strengths: the complete equilibrium set; \st{analytical}]\label{prop:middle}
Let $1/2<\tau(r)<1$, $\rho\DT(r)<2k$ and $\lambda_E(r)\ne\lb$. The equilibria are exactly these:
\begin{enumerate}[label=(\Alph*),nosep]
\item[(Z)] \emph{Zero research:} $\lh=0$, $\sigma_H=\sigma_L=\delta_0$, $\mu(q)=1/2$, $\zeta_q=0$; entry $\rho$.
  It always exists.
\item[(R)] \emph{The research branch:} $\lh=\lambda_R(r)$, $\sigma_H=\delta_1$, $\sigma_L=\delta_0$,
  $(\mu(-1),\mu(0),\mu(1))=(1/2,\ \hat b_0/(2\hat b_0+\lh),\ M_{\lh})$, $(\zeta_{-1},\zeta_0,\zeta_1)=(0,0,1)$. It exists
  if and only if $\varphi_a(r)\ge0$ and $F_r(\lb)>0$.
\item[($\bar{\rm R}$)] \emph{Research at the cap:} $\lh=\lb$ with the orders and entry of (R). It exists if and only if
  $F_r(\lb)\le0$ and $\lambda_E(r)\le\lb$.
\item[(T)] \emph{Threshold, rule CM only:} $\lh=\lambda_E(r)$, $\sigma_H=\delta_1$, $\sigma_L=\delta_0$, $\mu(1)=\tau(r)$,
  $\zeta_1=(e^*-\rho)/(1-\rho)$, $\zeta_0=\zeta_{-1}=0$. It exists if and only if $\lambda_E(r)<\lb$ and
  $\varphi_a(r)>0$.
\end{enumerate}
Every trade choice in these equilibria is strict, so the investor tie rule plays no role. Entry is $\rho$ in (Z),
$\rho+(1-\rho)[\beta+(1/2-\beta)\lh]$ in (R) and ($\bar{\rm R}$), and $\rho+(e^*-\rho)[\beta+(1/2-\beta)\lambda_E]$ in
(T).

At calibration B, at each $r\in\{31/20,8/5,33/20\}$: $\varphi_a(r)>0$, $F_r(\lb)>0$ and $\lambda_E(r)<\lb$. So the
equilibrium set is $\{{\rm Z},{\rm R}\}$ under rule C1 and $\{{\rm Z},{\rm R},{\rm T}\}$ under rule CM, under either
investor rule. Table~\ref{tab:middle} gives the values. \st{Computer-assisted} enclosures.
\end{proposition}

\begin{table}[t]
\centering\small
\caption{The intermediate strengths at calibration B ($a=\reg{res_a}$, $\lb=\reg{res_cap}$).}\label{tab:middle}
\begin{tabular}{lccc}
\toprule
 & $r=1.55$ & $r=1.60$ & $r=1.65$\\
\midrule
spread $\DT(r)$ & \reg{res_Delta_r155} & \reg{res_Delta_r160} & \reg{res_Delta_r165}\\
threshold $\tau(r)$ & \reg{res_tau_r155} & \reg{res_tau_r160} & \reg{res_tau_r165}\\
sale loss $k-\rho\DT/2$ & \reg{res_ULgap_r155} & \reg{res_ULgap_r160} & \reg{res_ULgap_r165}\\
recruitment intensity $\lambda_E(r)$ & \reg{res_lamE_r155} & \reg{res_lamE_r160} & \reg{res_lamE_r165}\\
research $\lambda_R(r)$ & \reg{res_lambda_r155} & \reg{res_lambda_r160} & \reg{res_lambda_r165}\\
$\lambda_R-\lambda_E$ & \reg{res_lamgap_r155} & \reg{res_lamgap_r160} & \reg{res_lamgap_r165}\\
$\lb-\lambda_R$ & \reg{res_capgap_r155} & \reg{res_capgap_r160} & \reg{res_capgap_r165}\\
purchase posterior $\mu(1)=M_{\lambda_R}$ & \reg{res_muplus_r155} & \reg{res_muplus_r160} & \reg{res_muplus_r165}\\
purchase payoff $U_H(1)=2a\lambda_R$ & \reg{res_UH_r155} & \reg{res_UH_r160} & \reg{res_UH_r165}\\
entry on the branch $\Ent_R$ & \reg{res_entry_r155} & \reg{res_entry_r160} & \reg{res_entry_r165}\\
high-class ownership $\Own_H$ & \reg{res_own_r155} & \reg{res_own_r160} & \reg{res_own_r165}\\
\midrule
threshold entry after a purchase $e^*$ (CM) & \reg{res_th_estar_r155} & \reg{res_th_estar_r160} & \reg{res_th_estar_r165}\\
entry in the threshold equilibrium (CM) & \reg{res_th_entry_r155} & \reg{res_th_entry_r160} & \reg{res_th_entry_r165}\\
\bottomrule
\end{tabular}

\smallskip
\parbox{0.92\linewidth}{\footnotesize Zero research has entry $\rho=0.25$ at each strength. $\DT$, $\tau$,
$\lambda_E$ and $e^*$ are exact rationals; the run log prints them. $\lambda_R$ is enclosed by exact rational
bisection of $F_r$ to width $10^{-24}$. The decimals are rounded from the enclosures.}
\end{table}

\begin{proposition}[The research branch; \st{analytical}; the interval at B is \st{computer-assisted}]\label{prop:branch}
\begin{enumerate}[label=(\alph*),nosep]
\item On $\{r:\DT(r)>2k\}$, the root $\lambda_R(r)$ is continuously differentiable and strictly rising:
\begin{equation}\label{eq:dlam}
\lambda_R'(r)=\frac{\beta\,[1-\lambda_R(r)]\,\DT'(r)}{\partial_\lambda F_r(\lambda_R(r))}>0 .
\end{equation}
On the research branch (R), entry $\Ent_R(r)=\rho+(1-\rho)[\beta+(1/2-\beta)\lambda_R(r)]$ and ownership
$\Own_H(r)=\tfrac12[\rho+(1-\rho)(\beta+(1-\beta)\lambda_R(r))]$ rise strictly with $r$.
\item At calibration B, (R) exists exactly for $r\in[r_R,\bar r_R)$, where $r_R=\reg{res_rR}$ is the unique root of
$\varphi_{1/40}$ on $[6/5,7/4]$ and $\bar r_R=\fr\left((2a\lb+k)(1+e)\right)=\reg{res_rcap}$, with
$\fr(D)=1+D+\sqrt{D^2+2D}$. Since $\bar r_R<r_N=\reg{res_rN}$, zero research is also an equilibrium at every
point of the branch. Along the branch, research rises from $\lambda_E(r_R)=\reg{res_lambda_rR}$ to the cap, and
entry rises from \reg{res_entry_rR} to \reg{res_entry_cap}.
\item At calibration B, research at the cap ($\bar{\rm R}$) coexists with zero research on $[\bar r_R,r_N)$, and
entry there equals the fixed-$\lambda$ level \reg{res_entry_cap}. Under rule CM the threshold equilibrium (T)
exists on $(r_R,r_N)$. Its entry falls across the three intermediate strengths
(Table~\ref{tab:middle}).
\end{enumerate}
\end{proposition}

\begin{proposition}[Strong incumbents; \st{analytical}, interval margins at B]\label{prop:strong}
Let $0<\beta\le1/3$, $0<\lb<1$, $\tau(r)>1/2$, $M_{\lb}\ne\tau(r)$, and
\begin{equation}\tag{U}\label{eq:U}
\rho\DT(r)\left(1-M_{\lb}\right)-k>a\lb .
\end{equation}
Then under every pair of tie rules the equilibrium is unique: $\lh=\lb$, $\sigma_H=\delta_1$, $\sigma_L=\delta_{-1}$,
$(\mu(-1),\mu(0),\mu(1))=(m_{\lb},1/2,M_{\lb})$, $\zeta_0=\zeta_{-1}=0$, and $\zeta_1=1$ if $M_{\lb}>\tau(r)$,
$\zeta_1=0$ if $M_{\lb}<\tau(r)$. Entry is $\rho+(1-\rho)(\bar b+\lb/2)$ with $\bar b=(1-\lb)\beta$ in the first case
and $\rho$ in the second. This is the fixed-$\lambda$ equilibrium with presence $\lb$.

Let instead $M_{\lb}=\tau(r)$, with the other hypotheses unchanged. Then every equilibrium has $\lh=\lb$,
$\sigma_H=\delta_1$, $\sigma_L=\delta_{-1}$ and $\zeta_0=\zeta_{-1}=0$. Under rule C1, $\zeta_1=1$ and the equilibrium
is unique. Under rule CM, each $\zeta_1\in[0,1]$ gives one equilibrium. Its entry
$\rho+(1-\rho)\zeta_1(\bar b+\lb/2)$ fills the interval $[\rho,\rho+(1-\rho)(\bar b+\lb/2)]$.

At calibration B, \eqref{eq:U} holds for every $r\in(r_U^R,h)$, with $r_U^R=\reg{res_rUR}$. On this interval
$\tau$ rises strictly, so $M_{\lambda_B}=\tau(r)$ holds only at the entry ceiling $r_C=\reg{res_rC}$. The
equilibrium is therefore unique at every $r\in(r_U^R,h)$ under rule C1. Under rule CM it is unique at every
$r\in(r_U^R,h)$ with $r\neq r_C$. At $r_C$ under rule CM, research is at the cap and entry fills
$[\rho,\reg{res_entry_cap}]$. Also $\tau(3)<M_{\lambda_B}=e/(1+e)<\tau(18/5)$. So at $r_1=3$ the unique outcome
has research at the cap and entry \reg{res_entry_r1}. At $r_2=18/5$ the unique outcome has research at the cap and
entry \reg{res_entry_r2}.
\end{proposition}

\begin{proposition}[The admissible region at B; \st{computer-assisted}]\label{prop:region}
Fix $\beta=0.294655$ and the other primitives of calibration B, and vary the cap $\lb$ and the cost $a$. Call the
statements below \emph{the claims}:
\begin{enumerate}[label=(\roman*),nosep]
\item Proposition~\ref{prop:weak} at $6/5$;
\item at each $r\in\{31/20,8/5,33/20\}$, the equilibrium set $\{{\rm Z},{\rm R}\}$ under C1 and
  $\{{\rm Z},{\rm R},{\rm T}\}$ under CM, with $\lambda_E(r)<\lambda_R(r)<\lb$;
\item at $3$ and $18/5$, the unique equilibrium of Proposition~\ref{prop:strong}, with research at the cap, entry
  above $\rho$ at $3$ and entry $\rho$ at $18/5$.
\end{enumerate}
The claims hold with strict margins if and only if
\begin{equation}\label{eq:region}
\lambda_E(3)<\lb<\lambda_E(18/5)\qquad\text{and}\qquad a_{\rm lo}(\lb)<a<a_{\rm up},
\end{equation}
where
\[
a_{\rm lo}(\lb)=\frac{\DT(33/20)\,m_{\lb}-k}{2\lb},\qquad
a_{\rm up}=\frac{\DT(31/20)\,[1-\tau(31/20)]-k}{2\lambda_E(31/20)}=\reg{res_a_up} .
\]
The strip is $(\reg{res_cap_lo},\reg{res_cap_hi})$. Table~\ref{tab:region} gives the cost window across the strip.
At the cap $\lambda_B$ the window is $(\reg{res_a_lo},\reg{res_a_up})$. At $a=1/40$ the cap range is
$(\reg{res_cap_lo_a},\reg{res_cap_hi})$.
\end{proposition}

\begin{corollary}[General $\beta$; \st{computer-assisted}]\label{cor:beta}
Keep $1+u=e$, set the cap to $\lb=\lambda_B(\beta)=u\beta/(1+u\beta)$, keep $a=1/40$, and keep the other primitives
of calibration B. Then the claims hold with strict margins for every $\beta\in(\beta_{\rm lo},1/3]$, with
$\beta_{\rm lo}=\reg{res_beta_lo}$. For $\beta\le\beta_{\rm lo}$, research reaches the cap at $33/20$.
\end{corollary}

\begin{proposition}[Price access at $r_1$; \st{analytical}, interval values at B]\label{prop:access}
Let $r=3$, $B_3(1/2)<c_H<B_3(M_{\lb})$, and let \eqref{eq:U} hold at $r=3$. Compare the equilibrium with price
access and the equilibrium of the economy in which the challenger does not see the price. Both economies have the
unique outcome with research at the cap and full trades. Access raises expected proceeds and expected acquisition
surplus net of preparation costs by
\begin{equation}\label{eq:access}
\Delta\mathcal R_T=(1-\rho)\left(\bar b+\tfrac{\lb}2\right)\left[t_L-t_0+\DT M_{\lb}\right],\qquad
\Delta\mathcal W=(1-\rho)\left(\bar b+\tfrac{\lb}2\right)\left[B_3(M_{\lb})-c_H+\frac{p^2}{r}\right].
\end{equation}
Both are positive. The research cost $(a/2)\lb^2$ and the trading cost $k\lb$ are the same in both economies. At
calibration B: proceeds rise from \reg{res_R_hidden_r1} to \reg{res_R_access_r1} (gain \reg{res_dR_r1}), and net
surplus rises from \reg{res_W_hidden_r1} to \reg{res_W_access_r1} (gain \reg{res_dW_r1}). The research cost is
\reg{res_cost_r1} and the trading cost is \reg{res_tcost_r1} in both economies.
\end{proposition}

\begin{remark}[What the quadratic cost does]\label{rem:cost}
The quadratic cost enters the proofs in three places.
\begin{enumerate}[nosep]
\item Strict concavity of $zV-(a/2)z^2$ gives a unique research response $\min\{\lb,V/a\}$
  (Lemma~\ref{lem:joint}). No equilibrium needs indifference over research.
\item The marginal cost is zero at zero research. So research is positive whenever $V>0$, and zero research is an
  equilibrium exactly when $V=0$ at the uninformative price, that is when $\rho\DT\le2k$. This condition does not
  involve $a$.
\item The marginal cost $a\lambda$ rises, and the research value $\frac12(\DT m_\lambda-k)$ on the branch falls,
  because more research removes camouflage. They meet once, at the root $\lambda_R(r)$ of the quadratic $F_r$. The
  root moves smoothly with $r$ by \eqref{eq:dlam}, so entry inherits a strict slope.
\end{enumerate}
With a linear research cost $c\lambda$, the research response is $0$ or $\lb$ except at $V=c$. Research that moves
with $r$ would then need investor indifference over research, and entry would be flat wherever research sits
at $0$ or at $\lb$. In the endpoint results the cost enters only through the marginal cost at the cap, $a\lb$, in
condition~\eqref{eq:U}.
\end{remark}

\begin{remark}[Tie rules]\label{rem:ties}
The investor tie rule plays no role at the six strengths, because every trade choice is strict. The challenger tie
rule matters at the six strengths only for the threshold equilibrium (T), which exists only under rule CM. Among
strong incumbents it matters only at the entry ceiling $r_C$ (Proposition~\ref{prop:strong}). Research never has a
tie. The research branch, its strict rise and the uniqueness at $6/5$, $3$ and $18/5$ hold under all four pairs of
rules.
\end{remark}

\begin{remark}[Start of the branch]\label{rem:start}
In the fixed-$\lambda$ model, informative trade first appears at $r_A=\reg{res_rA}$, where $\DT(1-\tau)=k$. With
research, the purchase posterior at the threshold must also pay the marginal research cost there:
$\DT(1-\tau)=k+2a\lambda_E(r)$, that is $\varphi_a(r_R)=0$. So $r_R>r_A$ at every $a>0$. On the fixed-$\lambda$
path, the purchase rate adjusts through a mixture of the high class. On the research branch, the research intensity
adjusts, and the high class buys for sure.
\end{remark}

\section{Proofs}\label{sec:proofs}

\subsection{Proof of Proposition~\ref{prop:weak}}

\emph{The candidate.} Take the tuple of the statement. (E1): with $\lh=0$, $f_H=f_L=\eta$, so $\mu(q)=1/2$. (E2):
every order has posterior $1/2$, so the belief at every price is $1/2$. (E3): $\tau(r)>1/2$, so
$B_r(1/2)<c_H$ and $\zeta_q=0$ under both challenger rules. (E4): \eqref{eq:price} with $e(q)=\rho$ and
$\mu(q)=1/2$ gives the stated prices. (E5): $U_H(1)=\rho\DT/2-k<0$,
$U_L(-1)=\rho\DT/2-k<0$, and the wrong-signed orders earn at most $-k$ (Lemma~\ref{lem:bounds}(c)). So zero is the
strict best order of each class. (E6) holds by truthful bidding. (E7): $V=0$, so $\min\{\lb,V/a\}=0=\lh$.

\emph{Uniqueness.} Take any candidate. The gross return $q(\E[V_T\mid j,q]-P(q))$ of a nonzero order is at most
$e(q)\DT\max\{\mu(q),1-\mu(q)\}\le\DT<k$. So $U_j(q)<0$ for $q\neq0$ against every price function and entry rule.
Then $V=0$, Lemma~\ref{lem:joint} gives $\lh=0$, and (E5) gives $\sigma_H=\sigma_L=\delta_0$. The rest follows as
above. At $r_0=6/5$, $\DT=\reg{res_Delta_r0}$ and $k-\DT=\reg{margin_res_weak_B}$. \qed

\subsection{Proof of Proposition~\ref{prop:middle}}

Throughout, $1/2<\tau(r)<1$, $\rho\DT(r)<2k$ and $\lambda_E(r)\ne\lb$.

\emph{Step 1: (Z) is an equilibrium.} The verification is that of Proposition~\ref{prop:weak}. The only change is
the bound used in (E5): $U_H(1)=U_L(-1)=\rho\DT/2-k<0$ by hypothesis.

\emph{Step 2: (R).} Let $\lh=\lambda_R(r)$, $\sigma_H=\delta_1$, $\sigma_L=\delta_0$. By \eqref{eq:masses},
$f_H=(\hat b,\hat b_0,\hat b+\lh)$ and $f_L=(\hat b,\hat b_0+\lh,\hat b)$ on $(-1,0,1)$. So
$\mu(-1)=1/2$, $\mu(0)=\hat b_0/(2\hat b_0+\lh)<1/2$ and $\mu(1)=M_{\lh}$. Suppose $\varphi_a(r)\ge0$ and $F_r(\lb)>0$.
Then $\DT(1-\tau)\ge k+2a\lambda_E>k$, and $1-\tau<1/2$ gives $\DT>2k$, so the root $\lambda_R(r)$ exists.
\begin{itemize}[nosep]
\item (E2)--(E4). The three posteriors differ, so by Lemma~\ref{lem:bounds}(b) the belief at $P(q)$ is $\mu(q)$.
  The prices follow from \eqref{eq:price}. $F_r$ rises, $F_r(\lambda_E)=-[2\beta+\lambda_E(1-2\beta)]\varphi_a(r)\le0=F_r(\lambda_R)$, so
  $\lambda_R\ge\lambda_E$ and $\mu(1)=M_{\lambda_R}\ge\tau(r)$. Rule C1 gives $\zeta_1=1$. Under rule CM, $\zeta_1=1$ is
  forced if $\mu(1)>\tau$ and allowed if $\mu(1)=\tau$. The posteriors $\mu(0),\mu(-1)\le1/2<\tau$ give
  $\zeta_0=\zeta_{-1}=0$.
\item (E5). $U_H(1)=\DT m_{\lh}-k$. The root condition $F_r(\lambda_R)=0$ gives $\DT m_{\lambda_R}-k=2a\lambda_R>0$.
  The wrong-signed orders earn at most $-k$. So the purchase is the strict best order of class $H$.
  $U_L(-1)=\rho\DT\mu(-1)-k=\rho\DT/2-k<0$, so zero is the strict best order of class $L$. Both investor rules agree.
\item (E7). $V=\frac12(2a\lambda_R)+\frac12\cdot0=a\lambda_R$. Since $F_r(\lb)>0$, $\lambda_R<\lb$. So
  $\min\{\lb,V/a\}=\lambda_R=\lh$.
\end{itemize}
Conversely, (R) needs $\mu(1)\ge\tau$, which is $\lambda_R\ge\lambda_E$, that is $\varphi_a\ge0$. It also needs
$\lambda_R<\lb$, that is $F_r(\lb)>0$.

\emph{Step 3: ($\bar{\rm R}$).} The orders and posteriors are those of Step 2 with $\lh=\lb$. If $\lambda_E\le\lb$, then
$\mu(1)=M_{\lb}\ge\tau$ and $\zeta_1=1$ as in Step 2. Then $V=\frac12(\DT m_{\lb}-k)$, and $F_r(\lb)\le0$ is
$V\ge a\lb$, so $\min\{\lb,V/a\}=\lb$. The purchase earns $\DT m_{\lb}-k\ge2a\lb>0$, so the order choices are strict as in
Step 2. The two conditions are also necessary.

\emph{Step 4: (T), rule CM.} Let $\lh=\lambda_E<\lb$, $\sigma_H=\delta_1$, $\sigma_L=\delta_0$, and
$\zeta_1=(e^*-\rho)/(1-\rho)$. Then $\mu(1)=M_{\lambda_E}=\tau(r)$, so $B_r(\mu(1))=c_H$ and rule CM allows any
$\zeta_1$. The definition \eqref{eq:phi} gives $U_H(1)=e^*\DT(1-\tau)-k=2a\lambda_E>0$, so the purchase is strict.
The low class abstains strictly as in Step 2. $V=a\lambda_E$, so $\min\{\lb,V/a\}=\lambda_E$. The preparation
probability lies in $(0,1)$ if and only if $\rho<e^*<1$. Here $e^*<1$ is $\varphi_a(r)>0$. Also
$e^*>\rho$ always, because $\rho\DT(1-\tau)<\rho\DT/2<k<2a\lambda_E+k$. At $e^*=1$ the tuple is (R) with
$\lambda_R=\lambda_E$.

\emph{Step 5: no other equilibrium.} Take any equilibrium.
\begin{itemize}[nosep]
\item If $\lh=0$: $f_j=\eta$, $\mu(q)=1/2$ and $e(q)=\rho$, so $U_H(1)=U_L(-1)=\rho\DT/2-k<0$. (E5) gives
  $\sigma_H=\sigma_L=\delta_0$. This is (Z).
\item If $\lh>0$: Lemma~\ref{lem:buy} gives $\sigma_H=\delta_1$, $\mu(1)=M_{\lh}$, $\mu(-1)\le1/2$ and $\zeta_{-1}=0$.
  Then $U_L(-1)=\rho\DT\mu(-1)-k\le\rho\DT/2-k<0$, and $\sigma_L=\delta_0$. So $V=\frac12(e(1)\DT m_{\lh}-k)^+$. Three
  cases remain.
  \begin{itemize}[nosep]
  \item $\lh<\lambda_E$: $\mu(1)<\tau$, $e(1)=\rho$, and $\rho\DT m_{\lh}<\rho\DT/2<k$, so $V=0$. Lemma~\ref{lem:joint} then
    gives $\lh=0$, a contradiction. This step excludes positive research with floor entry after a purchase.
  \item $\lh>\lambda_E$: $\mu(1)>\tau$, so $e(1)=1$ under both rules, and $V=\frac12(\DT m_{\lh}-k)^+$. If $\lh<\lb$, then
    $a\lh=V>0$, which is $F_r(\lh)=0$, so $\lh=\lambda_R$: this is (R). If $\lh=\lb$, then $V\ge a\lb$, which is
    $F_r(\lb)\le0$: this is ($\bar{\rm R}$).
  \item $\lh=\lambda_E$: under C1, $e(1)=1$ and the previous case applies with $\lambda_R=\lambda_E$. Under CM,
    $e(1)=e\in[\rho,1]$. Since $\lambda_E\ne\lb$ and $\lh=\lambda_E\le\lb$, research is interior, and
    $a\lambda_E=\frac12(e\DT(1-\tau)-k)^+$. So $e=e^*$: this is (T), or (R) when $e^*=1$.
  \end{itemize}
\end{itemize}
In every equilibrium of the list the purchase earns $2a\lh>0$, and the low-class sale earns $\rho\DT/2-k<0$. So
every trade choice is strict. Entry follows from \eqref{eq:entry}: the purchase has probability
$\frac12(\hat b+\lh)+\frac12\hat b=\beta+(1/2-\beta)\lh$, and only the purchase can raise entry above $\rho$.

\emph{Calibration B.} At the three strengths, Table~\ref{tab:middle} gives $k-\rho\DT/2>0$,
$\lambda_R-\lambda_E>0$ (equivalent to $\varphi_{1/40}>0$) and $\lb-\lambda_R>0$ (equivalent to $F_r(\lb)>0$). So
$\lambda_E<\lb$ too. Also $\rho<e^*<1$. The program enumerates the list above for each pair of tie rules, builds
each tuple, and recomputes (E1)--(E7) at the tuple. It prints every order payoff, every posterior and every price. \qed

\subsection{Proof of Proposition~\ref{prop:branch}}

(a) The polynomial $F_r(\lambda)$ is smooth in $(\lambda,r)$, since its coefficients are affine in $\DT(r)$. For
$\lambda\ge0$, $\partial_\lambda F_r=4a(1-2\beta)\lambda+4a\beta+k(1-2\beta)+\beta\DT>0$ and
$\partial_rF_r=-\beta(1-\lambda)\DT'(r)<0$, because $\lambda<1$ and $\DT'(r)=\frac12(1-\ell^2/r^2)>0$. The implicit
function theorem gives \eqref{eq:dlam}. The purchase has probability $\beta+(1/2-\beta)\lambda_R$, and only the
purchase raises entry, to $1$. So $\Ent_R=\rho+(1-\rho)[\beta+(1/2-\beta)\lambda_R]$. Since $\beta<1/2$, $\Ent_R$ rises
with $\lambda_R$ and hence with $r$. Class $H$ buys with probability $\hat b+\lambda_R=\beta+(1-\beta)\lambda_R$, which
gives $\Own_H$.

(b) The proof has four steps.
\begin{enumerate}[nosep]
\item \emph{$r\le6/5$.} $\DT(r)\le\DT(6/5)<k$, so Proposition~\ref{prop:weak} applies and (R) does not exist.
\item \emph{$\varphi_{1/40}$ rises on $J=[6/5,7/4]$.} By \eqref{eq:lamE},
  $\varphi_a'=\DT'(1-\tau)-\DT\tau'-2a\beta\tau'/\{\beta+(1-2\beta)(1-\tau)\}^2$. At the benchmark,
  $\tau'(r)=3(8r^2-r+2)/[2(20r-r^2-1)^2]>0$. On $J$ these bounds hold, since $\DT'$, $\DT$ and $\tau$ rise
  and $20r-r^2-1$ rises:
  $\DT'\ge\reg{res_bound_dDelta}$, $1-\tau\ge1-\tau(7/4)=\reg{res_bound_tau74}$, $\DT\le\DT(7/4)=\reg{res_bound_Delta74}$,
  and $\tau'\le\tau'_{\max}=\reg{res_bound_taup}$. These give
  $(\DT(1-\tau))'\ge\reg{res_dphi_lead}$ and
  $\varphi_a'\ge\reg{res_dphi_lead}-2a\beta\tau'_{\max}/\{\beta+(1-2\beta)(1-\tau(7/4))\}^2$.
  At $a=1/40$ the bound is $\varphi_a'\ge\reg{res_phip_lb}>0$. It stays positive for every
  $a<a_{\rm mono}=\reg{res_amono}$, and $a_{\rm mono}>a_{\rm up}$. Exact rational evaluation gives
  $\varphi_{1/40}(6/5)<0<\varphi_{1/40}(7/4)$. So $\varphi_{1/40}$ has the unique root $r_R$ on $J$, and
  bisection with exact signs encloses it.
\item \emph{$r_R\le r<\bar r_R$.} Here $r<r_N$, so $\rho\DT<2k$, and $1/2<\tau<1$. Also $\varphi_{1/40}(r)\ge0$.
  $\lambda_R$ rises by (a), and $\lambda_R(\bar r_R)=\lb$ by the definition of $\bar r_R$: with $m_{\lb}=1/(1+e)$,
  $F_r(\lb)=0$ is $\DT(r)=(2a\lb+k)(1+e)$, and $\fr$ inverts $\DT$ on $r>1$. So $F_r(\lb)>0$ for $r<\bar r_R$, and (R)
  exists by Proposition~\ref{prop:middle}. On $(6/5,r_R)$, $\varphi_{1/40}<0$ and (R) does not exist.
\item \emph{$r\ge\bar r_R$.} Take any equilibrium with $e(1)=1$ and $\lh<\lb$. Lemma~\ref{lem:buy} gives
  $\sigma_H=\delta_1$ and $\mu(1)=M_{\lh}$. The low class earns at least zero, so
  $a\lh=V\ge\frac12(\DT m_{\lh}-k)$, that is $F_r(\lh)\ge0$. Hence $\lh\ge\lambda_R(r)\ge\lambda_R(\bar r_R)=\lb$, a
  contradiction. So no equilibrium with a recruiting purchase and research below the cap exists for $r\ge\bar r_R$. This step
  covers $r>r_N$ too, where the low class may sell.
\end{enumerate}
Zero research is an equilibrium exactly when $\rho\DT\le2k$, that is $r\le r_N$, and $\bar r_R<r_N$. The endpoint
values are $\lambda_R(r_R)=\lambda_E(r_R)$ and $\Ent_R(\bar r_R)=\rho+(1-\rho)[\beta+(1/2-\beta)\lb]$.

(c) On $[\bar r_R,r_N)$, $F_r(\lb)\le0$ and $\lambda_E(r)\le\lambda_E(7/4)=\reg{res_lamE_r175}<\lb$, so ($\bar{\rm R}$) exists by
Proposition~\ref{prop:middle}. Entry is the expression of (a) at $\lambda=\lb$. On $(r_R,r_N)$ under rule CM,
$\varphi_{1/40}>0$ and $\lambda_E<\lb$, so (T) exists. Table~\ref{tab:middle} gives its entry at the three
strengths. \qed

\subsection{Proof of Proposition~\ref{prop:strong}}

\emph{Uniqueness.} Take any equilibrium. By Lemma~\ref{lem:bounds}(a) every posterior lies in
$[m_{\lb},M_{\lb}]$, whatever the order distributions. Entry is at least $\rho$ after every price. Hence
\[
U_H(1)=e(1)\DT(1-\mu(1))-k\ge\rho\DT(1-M_{\lb})-k,\qquad U_L(-1)=e(-1)\DT\mu(-1)-k\ge\rho\DT m_{\lb}-k ,
\]
and both lower bounds equal $\rho\DT(1-M_{\lb})-k>a\lb\ge0$ by \eqref{eq:U}. So the purchase and the sale are the
strict best orders of the two classes, under both investor rules. Then $V=\frac12[U_H(1)+U_L(-1)]>a\lb$, and
Lemma~\ref{lem:joint} gives $\lh=\lb$. With full orders, $f_H=(\bar b,\bar b_0,\bar b+\lb)$ and
$f_L=(\bar b+\lb,\bar b_0,\bar b)$, so $(\mu(-1),\mu(0),\mu(1))=(m_{\lb},1/2,M_{\lb})$. Since $\tau(r)>1/2$,
$\zeta_0=\zeta_{-1}=0$. Since $M_{\lb}\neq\tau(r)$, $\zeta_1$ is $1$ or $0$ under both challenger rules. Every
component of the tuple is therefore pinned down. Only this last step uses $M_{\lb}\neq\tau(r)$.

\emph{Existence.} The tuple of the statement satisfies (E1)--(E4) and (E6) by construction and
Lemma~\ref{lem:bounds}(b), (E5) by the strict bounds above, and (E7) by $V>a\lb$. The order law, the posteriors, the prices and the entry rule
equal those of the fixed-$\lambda$ equilibrium with presence $\lb$ and full orders.

\emph{Entry.} The purchase has probability $\frac12(\bar b+\lb)+\frac12\bar b=\bar b+\lb/2$. It recruits exactly
when $M_{\lb}>\tau(r)$.

\emph{The case $M_{\lb}=\tau(r)$.} The uniqueness argument above gives $\lh=\lb$, full orders,
$\mu(1)=M_{\lb}=\tau(r)$ and $\zeta_0=\zeta_{-1}=0$. At $\mu(1)=\tau(r)$, $B_r(\mu(1))=c_H$, so the challenger with the
high cost is indifferent after a purchase. Rule C1 sets $\zeta_1=1$, and the tuple is unique. Rule CM permits each
$\zeta_1\in[0,1]$. Fix one, set $e(1)=\rho+(1-\rho)\zeta_1$, and set $P(1)$ by \eqref{eq:price}. The existence step
holds for this tuple without change: (E5) and (E7) use only the bounds above, and these bounds use only
$e(q)\ge\rho$. Since $\mu(1)>1/2\ge\mu(0),\mu(-1)$ and $e(1)\ge\rho=e(0)=e(-1)$, the purchase price differs from the
other two prices, so (E2) holds. Entry is $\rho+(1-\rho)\zeta_1(\bar b+\lb/2)$.

\emph{Calibration B.} At $\lb=\lambda_B$, $M_{\lb}=e/(1+e)$. Condition~\eqref{eq:U} is
$\DT(r)>(k+a\lb)(1+e)/\rho$, which holds for $r>r_U^R=\fr((k+a\lb)(1+e)/\rho)$. Its margin is
\reg{margin_res_bound_r1_B} at $r=3$ and \reg{margin_res_bound_r2_B} at $r=18/5$. The posterior margins are
$M_{\lb}-\tau(3)=\reg{margin_res_tau_r1_B}$ and $\tau(18/5)-M_{\lb}=\reg{margin_res_tau_r2_B}$. The entry ceiling
$r_C$ solves $B_r(M_{\lb})=c_H$. The closed forms hold for $\ell<r<h$, and there $\DT'(r)>0$, so \eqref{eq:U}
holds on all of $(r_U^R,h)$. The numerator $c_H-g_L$ of $\tau$ rises in $r$, because $g_L=(\ell^2-p^2)/(2r)$ falls.
The denominator $g_H-g_L=h-r/2-\ell^2/(2r)$ is positive on $(\ell,h)$ and falls, because its derivative is
$-\frac12+\ell^2/(2r^2)<0$. So $\tau$ rises strictly, and $r_C$ is the only root of $M_{\lb}=\tau(r)$ on $(r_U^R,h)$.
The run log prints the margin of \eqref{eq:U} at $r_C$. \qed

\subsection{Proof of Proposition~\ref{prop:region}}

\emph{Sufficiency.} Let \eqref{eq:region} hold.
\begin{enumerate}[nosep]
\item Claim (i) holds for every $(\lb,a)$, since $\tau(6/5)>1/2$ and $\DT(6/5)<k$.
\item Claim (ii). The hypotheses $1/2<\tau<1$ and $\rho\DT<2k$ of Proposition~\ref{prop:middle} do not involve
  $(\lb,a)$ and hold at the three strengths. The condition $a<a_{\rm up}$ is $\varphi_a(31/20)>0$. Since
  $a<a_{\rm up}<a_{\rm mono}$, $\varphi_a$ rises on $J$, so $\varphi_a>0$ at $8/5$ and $33/20$ too. The condition
  $a>a_{\rm lo}(\lb)$ is $F_{33/20}(\lb)>0$, that is $\lambda_R(33/20)<\lb$. Since $\lambda_R$ rises in $r$,
  $\lambda_R<\lb$ at $31/20$ and $8/5$ too. Then $\lambda_E<\lambda_R<\lb$, and Proposition~\ref{prop:middle} gives the
  sets.
\item Claim (iii). On the strip, $\tau(3)<M_{\lb}<\tau(18/5)$, since $M_\lambda$ rises in $\lambda$ and
  $M_{\lambda_E(r)}=\tau(r)$. Condition~\eqref{eq:U} at $r=3$ is $a<a_{U3}(\lb):=[\rho\DT(3)m_{\lb}-k]/\lb$, and at
  $r=18/5$ it is $a<a_{U36}(\lb):=[\rho\DT(18/5)m_{\lb}-k]/\lb$. Both bounds fall in $\lb$, so on the strip they
  exceed their values at $\lb=\lambda_E(18/5)$, which are \reg{res_region_aU3_6} and \reg{res_region_a36_6}. Both
  exceed $a_{\rm up}$. So Proposition~\ref{prop:strong} applies at both strengths.
\end{enumerate}

\emph{Necessity.} Each failure of \eqref{eq:region} breaks one claim.
\begin{itemize}[nosep]
\item $\lb\le\lambda_E(3)$: every posterior at $3$ is at most $M_{\lb}\le\tau(3)$ (Lemma~\ref{lem:bounds}(a)). Entry
  above $\rho$ then needs a challenger tie, so claim (iii) fails with strict margins.
\item $\lb\ge\lambda_E(18/5)$: research at the cap gives $\mu(1)=M_{\lb}\ge\tau(18/5)$, so rule C1 raises entry above
  $\rho$, and claim (iii) fails.
\item $a\ge a_{\rm up}$: $\varphi_a(31/20)\le0$, so (R) is absent at $31/20$, or has $\mu(1)=\tau$ with a challenger
  tie.
\item $a\le a_{\rm lo}(\lb)$: $F_{33/20}(\lb)\le0$, so research at $33/20$ sits at the cap.
\end{itemize}
The window is nonempty on the whole strip: $a_{\rm lo}$ falls in $\lb$, because its numerator
$\DT(33/20)m_{\lb}-k$ falls and stays positive there, and $a_{\rm lo}(\lambda_E(3))=\reg{res_region_alo_0}<a_{\rm up}$.
The program evaluates every bound in exact rational arithmetic, or in interval arithmetic where $e$ enters. \qed

\subsection{Proof of Corollary~\ref{cor:beta}}

At $\lb=\lambda_B(\beta)$, $u_{\lb}=u=e-1$ for every $\beta$, so $M_{\lb}=e/(1+e)$ and $m_{\lb}=1/(1+e)$ do not
depend on $\beta$. Hence the endpoint conditions $\tau(3)<M_{\lb}<\tau(18/5)$ hold for every $\beta$.
Condition~\eqref{eq:U} at $3$ holds for every $\beta\le1/3$, because $\lambda_B(\beta)$ rises in $\beta$ and
$\rho\DT(3)/(1+e)-k-a\lambda_B(1/3)>0$. At the intermediate strengths,
$a_{\rm lo}=[\DT(33/20)/(1+e)-k]/[2\lambda_B(\beta)]$ falls in $\beta$. So $a_{\rm lo}<1/40$ exactly when
$\beta>\beta_{\rm lo}$. The upper bound $a_{\rm up}(\beta)=[\DT(31/20)(1-\tau)-k]/[2\lambda_E(31/20;\beta)]$ exceeds
$1/40$ for every $\beta<\reg{res_beta_hi}$. Indeed, by \eqref{eq:lamE},
$\partial\lambda_E/\partial\beta=(2\tau-1)(1-\tau)/\{\beta+(1-2\beta)(1-\tau)\}^2>0$, so $a_{\rm up}(\beta)$ falls in
$\beta$, and $\beta_{\rm hi}=\reg{res_beta_hi}$ is the point where it equals $1/40$. The monotonicity bound of
Proposition~\ref{prop:branch}(b) is $a_{\rm mono}(\beta)=L\,D(\beta)^2/(2\beta\tau'_{\max})$, where $L$ is the bound on
$(\DT(1-\tau))'$ and $D(\beta)=c+\beta(1-2c)$ with $c=1-\tau(7/4)=\reg{res_bound_tau74}$. The derivative of
$D(\beta)^2/\beta$ has the sign of $\beta(1-2c)-c$, which is negative for $\beta\le1/3$ because
$(1-2c)/3<c$. So $a_{\rm mono}$ falls in $\beta$, and at $\beta=1/3$ it equals $a_{\rm mono}=\reg{res_amono_third}>1/40$. The steps of Proposition~\ref{prop:region} then
apply at each $\beta\in(\beta_{\rm lo},1/3]$. \qed

\subsection{Proof of Proposition~\ref{prop:access}}

\emph{The hidden-price economy.} The challenger's belief stays at $1/2$. Since $B_3(1/2)<c_H$, $\zeta_q=0$ and
$e(q)=\rho$ after every order. The bounds of the proof of Proposition~\ref{prop:strong} hold with $e(q)=\rho$, so
\eqref{eq:U} again gives strict full orders and research at the cap, uniquely. With access,
Proposition~\ref{prop:strong} gives the same orders and research, and $\zeta_1=1$ because $M_{\lb}>\tau(3)$.

\emph{Costs.} The order law is the same in both economies. So the research cost $(a/2)\lb^2$ and the expected trading
cost $\lb k$ coincide, and these costs cancel from the comparison also when they use resources.

\emph{Proceeds.} Expected proceeds are $\sum_j\frac12\sum_qf_j(q)[t_0+e(q)(t_j-t_0)]$. Entry differs only after the
purchase, where it is $1$ instead of $\rho$. So
$\Delta\mathcal R_T=(1-\rho)\frac12[(\bar b+\lb)(t_H-t_0)+\bar b(t_L-t_0)]=(1-\rho)(\bar b+\lb/2)[t_L-t_0+\DT M_{\lb}]$,
because $t_H=t_L+\DT$ and $M_{\lb}=(\bar b+\lb)/(2\bar b+\lb)$. Every factor is positive.

\emph{Net surplus.} Write $S_0=\E[R\,\mathbf 1\{R\ge p\}]$ and $S_j=\E[\max\{\theta_j,R\}]$. For $\theta\ge p$,
$\max\{\theta,R\}-R\,\mathbf 1\{R\ge p\}=(\theta-\max\{p,R\})_++p\,\mathbf 1\{R<p\}$. In expectation, a prepared
challenger of class $j$ adds $S_j-S_0=g_j+p\Pr(R<p)=g_j+p^2/r$. The added high-cost
preparation after a purchase costs $c_H$. So
$\Delta\mathcal W=(1-\rho)(\bar b+\lb/2)[B_3(M_{\lb})-c_H+p^2/r]$. It is positive because
$B_3(M_{\lb})>c_H$. The program also computes both levels by direct summation over orders and classes, and the
differences agree with \eqref{eq:access} exactly. \qed

\section{Margins and the admissible region}\label{sec:margins}

Every hypothesis is a strict inequality. Table~\ref{tab:margins} records its margin at calibration B. The
hypotheses define an open set of primitives: each is a strict inequality between continuous functions of
$(h,\ell,p,\rho,c_H,k,\beta,\lb,a,r)$. The admissible region of Proposition~\ref{prop:region} is the slice of this set
in $(\lb,a)$ at the other benchmark primitives.

\begin{table}[ht]
\centering\small
\caption{Strict hypotheses and their margins at calibration B.}\label{tab:margins}
\begin{tabular}{lll}
\toprule
Hypothesis & Used in & Margin at B\\
\midrule
$B_r(1/2)<c_H$ (so $\tau>1/2$), at $r=6/5$ (smallest of six) & Lemma~\ref{lem:buy}, Prop.~\ref{prop:weak} & \reg{margin_res_half_B}\\
$\DT(6/5)<k$ & Prop.~\ref{prop:weak} & \reg{margin_res_weak_B}\\
$\rho\DT(33/20)<2k$ (largest middle strength) & Prop.~\ref{prop:middle} & \reg{margin_res_nosale_B}\\
$\varphi_{1/40}(31/20)>0$ & Prop.~\ref{prop:middle} & \reg{res_phi_r155}\\
\quad equivalently $a<a_{\rm up}$ & & \reg{margin_res_aup_B}\\
$\lambda_R(33/20)<\lb$ & Prop.~\ref{prop:middle} & \reg{margin_res_lamcap_B}\\
\quad equivalently $a>a_{\rm lo}(\lb)$ & & \reg{margin_res_alo_B}\\
$\lambda_E(3)<\lb$ & Prop.~\ref{prop:strong} & \reg{margin_res_caplo_B}\\
\quad equivalently $\tau(3)<M_{\lb}$ & & \reg{margin_res_tau_r1_B}\\
$\lb<\lambda_E(18/5)$ & Prop.~\ref{prop:strong} & \reg{margin_res_caphi_B}\\
\quad equivalently $M_{\lb}<\tau(18/5)$ & & \reg{margin_res_tau_r2_B}\\
\eqref{eq:U} at $r=3$ & Props.~\ref{prop:strong}, \ref{prop:access} & \reg{margin_res_bound_r1_B}\\
\eqref{eq:U} at $r=18/5$ & Prop.~\ref{prop:strong} & \reg{margin_res_bound_r2_B}\\
$\beta<1/3$ (the domain needs only $\beta\le1/3$) & Lemma~\ref{lem:bounds} & \reg{margin_res_beta_B}\\
$r_R<31/20$ & Prop.~\ref{prop:branch} & \reg{margin_res_rR_B}\\
$33/20<\bar r_R$ & Prop.~\ref{prop:branch} & \reg{margin_res_rcap_B}\\
$\bar r_R<r_N$ & Prop.~\ref{prop:branch} & \reg{margin_res_rN_B}\\
$a<a_{\rm mono}$ & Prop.~\ref{prop:branch} & \reg{margin_res_amono_B}\\
\bottomrule
\end{tabular}
\end{table}

The tightest hypothesis is $\lb<\lambda_E(18/5)$ at the very strong incumbent. The table gives its margin in two
coordinates:
\begin{itemize}[nosep]
\item in research intensity, $\lambda_E(18/5)-\lb=\reg{margin_res_caphi_B}$;
\item in posterior probability, $\tau(18/5)-M_{\lb}=\reg{margin_res_tau_r2_B}$.
\end{itemize}
The two numbers measure the same boundary, and they are not equal. They have the same sign because $M_\lambda$
rises strictly in $\lambda$ and $M_{\lambda_E(r)}=\tau(r)$. The margin comes from the choice of the cap,
$\lb=\lambda_B$, and so from the fixed-$\lambda$ calibration. Research does not change it, because research sits at
the cap at $18/5$. The margins of the table have different units, so the table does not rank them.

\begin{table}[ht]
\centering\small
\caption{The admissible region at calibration B: the cost window across the cap strip.}\label{tab:region}
\begin{tabular}{lccccc}
\toprule
Cap $\lb$ & $a_{\rm lo}(\lb)$ & $a_{\rm up}$ & window width & $a_{U3}(\lb)$ & $a_{U36}(\lb)$\\
\midrule
$\lambda_E(3)=\reg{res_region_cap_0}$ & \reg{res_region_alo_0} & \reg{res_a_up} & \reg{res_region_width_0} & \reg{res_region_aU3_0} & \reg{res_region_a36_0}\\
\reg{res_region_cap_1} & \reg{res_region_alo_1} & \reg{res_a_up} & \reg{res_region_width_1} & \reg{res_region_aU3_1} & \reg{res_region_a36_1}\\
\reg{res_region_cap_2} & \reg{res_region_alo_2} & \reg{res_a_up} & \reg{res_region_width_2} & \reg{res_region_aU3_2} & \reg{res_region_a36_2}\\
\reg{res_region_cap_3} & \reg{res_region_alo_3} & \reg{res_a_up} & \reg{res_region_width_3} & \reg{res_region_aU3_3} & \reg{res_region_a36_3}\\
\reg{res_region_cap_4} & \reg{res_region_alo_4} & \reg{res_a_up} & \reg{res_region_width_4} & \reg{res_region_aU3_4} & \reg{res_region_a36_4}\\
$\lambda_B=\reg{res_region_cap_5}$ & \reg{res_region_alo_5} & \reg{res_a_up} & \reg{res_region_width_5} & \reg{res_region_aU3_5} & \reg{res_region_a36_5}\\
$\lambda_E(18/5)=\reg{res_region_cap_6}$ & \reg{res_region_alo_6} & \reg{res_a_up} & \reg{res_region_width_6} & \reg{res_region_aU3_6} & \reg{res_region_a36_6}\\
\bottomrule
\end{tabular}

\smallskip
\parbox{0.92\linewidth}{\footnotesize The claims hold for $\lambda_E(3)<\lb<\lambda_E(18/5)$ and
$a_{\rm lo}(\lb)<a<a_{\rm up}$. The end rows are the open ends of the strip. The columns $a_{U3}$ and $a_{U36}$ are the
cost limits of condition~\eqref{eq:U} at $3$ and $18/5$; they exceed $a_{\rm up}$ on the whole strip. The cost
$a=1/40$ lies in the window for every cap above \reg{res_cap_lo_a}.}
\end{table}

\section{Certificates}\label{sec:cert}

\paragraph{Exact program.} \texttt{scripts/P3\_research/P3\_research.py}, run log
\texttt{scripts/P3\_research/P3\_research.log}. The program uses \texttt{fractions.Fraction} for every rational
quantity and a rational interval class for every quantity that depends on $e$. The enclosure of $e$ is the partial
sum $\sum_{n\le40}1/n!$ plus the remainder bound $2/41!$; its width is below $10^{-49}$. Roots of $F_r$ and of
$\varphi_{1/40}$ come from bisection with exact signs; square roots come from rational bisection. The program
builds the tuple of every equilibrium of Propositions~\ref{prop:weak}--\ref{prop:strong} for each pair of tie rules.
It then recomputes (E1)--(E7) at the tuple. It prints the posteriors, entry, prices, the payoff of every order of
both classes, the research value and the research best response. The program also reproduces the earlier values at calibration A
($\beta=1/3$, $\lb=6/17$): research \reg{res_calA_lambda_r155}, \reg{res_calA_lambda_r160}, \reg{res_calA_lambda_r165}
and entry \reg{res_calA_entry_r155}, \reg{res_calA_entry_r160}, \reg{res_calA_entry_r165}.
It writes \texttt{registry/P3\_research.csv} and the macro file that this document reads.

\paragraph{Printed enclosures.} Every enclosure in the run log and in the registry notes comes directly from the
exact rational endpoints. The lower end rounds down and the upper end rounds up, so each printed interval contains
the exact interval. An exact value prints as one decimal if it terminates; otherwise it prints as an outward pair.
The file \texttt{scripts/P3\_research/P3\_enclosures.json} keeps the endpoints of every registry value. A short
endpoint is the exact rational. A long endpoint is an outward decimal with 60 places. The program
\texttt{P3\_compare\_round2.py} (log \texttt{P3\_compare\_round2.log}) compares each registry enclosure with the
independent exact enclosure in \texttt{scripts/XP3/exact\_enclosures.json}. Every registry enclosure contains it.

\paragraph{Diagnostic scan (\st{numerical diagnostic}).} \texttt{scripts/P3\_research/P3\_float\_scan.py}, run log
\texttt{P3\_float\_scan.log}. The scan does not use Lemma~\ref{lem:buy} or the classification of Proposition~\ref{prop:middle}. It evaluates every
grid point $(\lh,\sigma_H(1),\sigma_L(-1))$ under rule C1, with $4001\times51\times51$ points per strength. It keeps
the points with order regret at most $5\cdot10^{-4}$ and research residual at most $3\cdot10^{-4}$. At each of the six
strengths, the clusters of kept points are exactly the equilibria of
Propositions~\ref{prop:weak}--\ref{prop:strong} under rule C1. There is one cluster at $6/5$, $3$ and $18/5$. There are
two clusters at each intermediate strength: zero research and the branch.

\end{document}
